\documentclass[11pt]{article}
\usepackage[T1]{fontenc}
\usepackage{lmodern}
\usepackage[margin=1in]{geometry}
\usepackage{amsmath,amssymb,amsthm,mathtools}
\usepackage{microtype,booktabs,array}
\usepackage{xcolor}
\usepackage[colorlinks=true,linkcolor=blue!45!black,citecolor=blue!45!black,urlcolor=blue!45!black]{hyperref}
\hypersetup{pdftitle={Growth and irrationality in local redirection trees},pdfauthor={Christopher D. Long}}
\setlength{\emergencystretch}{2em}
\allowdisplaybreaks
\newtheorem{theorem}{Theorem}[section]
\newtheorem{proposition}[theorem]{Proposition}
\newtheorem{lemma}[theorem]{Lemma}
\newtheorem{corollary}[theorem]{Corollary}
\theoremstyle{definition}
\newtheorem{definition}[theorem]{Definition}
\theoremstyle{remark}
\newtheorem{remark}[theorem]{Remark}
\numberwithin{equation}{section}
\title{Growth and irrationality in local redirection trees}
\author{Christopher D. Long}
\date{Research draft --- November 23, 2026}
\begin{document}
\maketitle
\begin{abstract}
We study the local redirection family $\operatorname{VAN}(\omega)$ for
finite $\omega\ge0$, including the PAN and DAN models. An exact
Crump--Mode--Jagers representation identifies the exponent
$\alpha_\omega\in(0,1)$ as the unique root of an explicit Euler--Lotka
equation. For every finite tree seed with at least two vertices, the
number of nonleaf vertices divided by $n^{\alpha_\omega}$ converges
almost surely to a finite strictly positive random variable. The root
is effectively computable whenever $\omega$ is computable. We prove
that the exponent is irrational for every finite algebraic
$\omega\ge0$: at a rational candidate exponent the transform
$\int_0^1 x^s e^{x-1}\,dx$ is transcendental by the
Siegel--Shidlovskii theorem, whereas the characteristic equation
requires an algebraic value. The equation itself was derived by
Hartle, Krapivsky, Redner, and Zhang; the arguments here supply a
branching-process proof and an arithmetic deduction. These local
models differ from the original interacting random friend tree.
\end{abstract}

\noindent\textbf{Keywords.} Random trees; redirection; Crump--Mode--Jagers
processes; Malthusian parameter; $E$-functions; irrationality.

\medskip
\noindent\textbf{MSC 2020.} 60J80, 05C80, 11J81.

\section{Introduction}
\label{intro:section}

The random friend-tree rule attaches a new vertex to a uniformly chosen
neighbour of a uniformly chosen existing vertex. A basic question is
the growth of its \emph{core}, the set of vertices of degree at least
two. Krapivsky and Redner studied this model in their work on emergent
network modularity~\cite{KR2017}. Its transition rates couple the local
environments of adjacent core vertices.

Hartle, Krapivsky, Redner, and Zhang introduced local redirection models
in which a core can redirect only to a leaf~\cite{HRKZ2026}. A parameter
$\omega$ controls direct attachment to cores with no leaf neighbours.
The two cases $\omega=0$ and $\omega=1$ are respectively prohibited
attachment to the nucleus (PAN) and direct attachment to the nucleus
(DAN); the full family is denoted $\operatorname{VAN}(\omega)$. Their
analysis gives the characteristic equation
\begin{equation}\label{intro:characteristic}
 \int_0^1 x^s e^{x-1}\,dx
 =\frac{(1-s)(s+\omega)}{\omega+(1-\omega)s},
 \qquad 0<s<1.
\end{equation}
Equation~\eqref{intro:characteristic} is
\cite[Eq.~(25)]{HRKZ2026}, with the PAN and DAN specializations in
Eqs.~(22) and~(17) of that paper.

We identify this root with the almost-sure core-growth exponent by an
exact branching construction, and prove a finite strictly positive
random amplitude for every finite seed. The proof eliminates the
waiting time at zero in a core's leaf-count process, applies the CMJ
law of large numbers, and controls the graph-size clock by a leaf-count
martingale. Section~\ref{van:section} also gives a certified root
algorithm.

For the arithmetic conclusion, we prove that
$F(r)=\int_0^1 x^r e^{x-1}\,dx$ is transcendental for every rational
$r>-1$. An explicit common-denominator bound makes the relevant
entire function an $E$-function, to which the classical
Siegel--Shidlovskii theorem applies. Equation~\eqref{intro:characteristic}
then implies irrationality of $\alpha_\omega$ for every finite
algebraic $\omega\ge0$ (Theorem~\ref{vana:irrationality}).
Section~\ref{vana:transfer-section} records the corresponding
specialization nonvanishing criterion for more general reductions.

Throughout the paper, $\alpha_\omega$ denotes the exponent of the
local model. We reserve $\mu$ for the exponent of the original
isotropic-redirection friend tree. The independent-family construction
used here applies to $\alpha_\omega$ and supplies no arithmetic
classification of $\mu$.

\section{Model and exact growth theorem}
\label{van:section}

The local redirection models of Hartle, Krapivsky, Redner, and Zhang
redirect a core with positive leaf count to one of its leaves
\cite{HRKZ2026}. The resulting independence of different core families
is the basis of the proof below.

Fix a finite parameter $\omega\ge0$.  A \emph{core} is a vertex of degree
at least two, and its \emph{leaf count} is the number of its neighbours
of degree one.  In continuous time the rules are:
\begin{enumerate}
\item Each leaf has a rate-one clock.  When it rings, attach a new vertex
to that leaf's unique neighbour.
\item Each core with positive leaf count has a rate-one clock.  When it
rings, choose one of its leaves uniformly and attach a new vertex to
that leaf.  The chosen leaf becomes a new core.
\item Each core with zero leaves has a rate-$\omega$ clock.  When it
rings, attach a new vertex directly to that core.
\end{enumerate}
The embedded tree at successive vertex births is precisely
$\operatorname{VAN}(\omega)$: $\omega=0$ gives PAN, after deleting
rejected selections, and $\omega=1$ gives DAN.  These rules are understood
after the initial seed has at least three vertices.  A two-vertex seed
reaches a three-vertex path after its first event and is included by
this finite initial delay.

Write $C(t)$, $L(t)$ and $C_0(t)$ for the numbers of cores, leaves and
zero-leaf cores.  Then $N(t)=C(t)+L(t)$ and the total vertex-birth rate is
\[
 L(t)+C(t)-C_0(t)+\omega C_0(t)
 =N(t)+(\omega-1)C_0(t)\le \max\{1,\omega\}N(t).
\]
In particular, the process is nonexplosive.  Let $C_n$ be the core count
when the tree has $n$ vertices.

\begin{theorem}[Exact modified-model exponent and almost-sure amplitude]
\label{van:main}
For every finite $\omega\ge0$, define, for $s>0$,
\begin{align}
 F(s)&=\int_0^1 x^s e^{x-1}\,dx,
 \label{van:F}\\
 m_\omega(s)&=
 \frac{(s+\omega)(1-F(s))}
 {s\bigl(s+\omega-\omega F(s)\bigr)}.
 \label{van:m}
\end{align}
There is a unique $\alpha_\omega\in(0,1)$ satisfying
$m_\omega(\alpha_\omega)=1$.  Equivalently,
\begin{equation}
 \boxed{\quad
 F(\alpha_\omega)=
 \frac{(1-\alpha_\omega)(\alpha_\omega+\omega)}
 {\omega+(1-\omega)\alpha_\omega}.
 \quad}
 \label{van:root}
\end{equation}
For every finite tree seed with at least two vertices, there are random
variables $A,V\in(0,\infty)$ such that, almost surely,
\begin{equation}
 e^{-\alpha_\omega t}C(t)\longrightarrow A,
 \qquad e^{-t}N(t)\longrightarrow V,
 \qquad
 \frac{C_n}{n^{\alpha_\omega}}\longrightarrow A V^{-\alpha_\omega}.
 \label{van:limits}
\end{equation}
The exponent $\alpha_\omega$ is effectively computable whenever
$\omega$ is computable.  More quantitatively,
\begin{equation}
 -m_\omega'(s)\ge\frac1{(1+s)^2}\ge\frac14,
 \qquad 0<s<1.
 \label{van:slope}
\end{equation}
\end{theorem}

The characteristic equation~\eqref{van:root} is the equation derived
in \cite[Eq.~(25)]{HRKZ2026}. The theorem identifies its unique root
with the almost-sure exponent and proves the positive-amplitude
limit for every finite seed. The parameters $\alpha_0$ and $\alpha_1$
consequently satisfy
$F(\alpha_0)=1-\alpha_0$ and $F(\alpha_1)=1-\alpha_1^2$.
Throughout the theorem, $\omega$ is finite.  Its assertion does not
include the separate limiting model $\operatorname{VAN}(\infty)$.

\subsection{The exact life process of a core}

Regard a newly created core as an individual whose children are the
cores it subsequently creates.  Every such newborn individual starts
with one leaf.  Its own leaf count $X$ has transition rates
\begin{equation}
 \ell\longrightarrow\ell+1 \text{ at rate }\ell,\qquad
 \ell\longrightarrow\ell-1 \text{ at rate }1
 \quad(\ell\ge1),\qquad
 0\longrightarrow1 \text{ at rate }\omega.
 \label{van:life}
\end{equation}
Each downward transition creates one child, itself starting at leaf
count one.  The parent loses the promoted leaf and the child acquires
the newly attached leaf.

The life processes of different individuals can be constructed
independently.  Their transition rates do not involve their core
neighbours; fresh future clock increments at a promotion are
independent of the past.  Thus this is a genuine single-type
Crump--Mode--Jagers process.  The exterior core edge of a newborn
individual does not affect its life process.  This observation also
allows us to study one genealogical family in isolation.

\begin{lemma}[An exact excursion transform]
\label{van:excursion}
For the chain \eqref{van:life} started at $\ell\ge1$, let $T_0$ be its
first hit of zero, prior to any restart from zero.  For $s>0$,
\begin{equation}
 h_\ell(s):=\mathbb E_\ell[e^{-sT_0};T_0<\infty]
 =\frac1{(\ell-1)!}\int_0^1
 x^s(1-x)^{\ell-1}e^{x-1}\,dx.
 \label{van:hitting}
\end{equation}
In particular, $h_1(s)=F(s)$.
\end{lemma}

\begin{proof}
Set $h_0=1$ and define $h_\ell$ by the integral on the right.
Integration of the derivative of
$x^{s+1}(1-x)^{\ell-1}e^{x-1}$ gives
\[
 (s+\ell+1)h_\ell=\ell h_{\ell+1}+h_{\ell-1},
 \qquad \ell\ge1.
\]
For $\ell=1$ the boundary value at $x=1$ is one and supplies the
term $h_0$; for $\ell\ge2$ both boundary values vanish.
Moreover $0\le h_\ell\le1/\ell!$.
The displayed recurrence is the discounted backward equation.
Apply its bounded stopped-martingale identity up to $T_0$, time $t$,
and the first hit of level $R$.  The contribution at level $R$
vanishes as $R\to\infty$, since $h_R\to0$; nonexplosion justifies
removing this localization.  The remaining time-$t$ term is bounded
by $e^{-st}$ and vanishes as $t\to\infty$.  This proves
\eqref{van:hitting}.
\end{proof}

Let $\xi$ be the reproduction point process of an individual started
at one, and let $\nu_\omega=\mathbb E\xi$ be its mean reproduction
measure.  The downward-event intensity gives
\begin{equation}
 \widehat\nu_\omega(s)
 =\mathbb E\int_0^\infty e^{-st}\xi(dt)
 =\mathbb E_1\int_0^\infty e^{-st}\mathbf1_{\{X_t>0\}}\,dt.
 \label{van:reproduction}
\end{equation}
Before $T_0$ this expected occupation is $(1-F(s))/s$.
Upon hitting zero, the chain waits an independent
$\operatorname{Exp}(\omega)$ time before restarting from one.
For $\omega=0$ this waiting time is infinite, with discounted
transform zero.  The strong Markov property therefore gives
\[
 \widehat\nu_\omega(s)
 =\frac{1-F(s)}s+
 F(s)\frac{\omega}{s+\omega}\widehat\nu_\omega(s).
\]
Solving proves $\widehat\nu_\omega(s)=m_\omega(s)$ as in
\eqref{van:m}; all its displayed denominators are positive for $s>0$.

The first child is born at age $B\sim\operatorname{Exp}(1)$.
Indeed, before the first core-clock ring only upward leaf moves
can occur, so the core remains active.  More generally, the entire
reproduction process is obtained by accepting a subset of the points
of a rate-one Poisson process $\Pi$, rejecting its points when $X=0$.
Consequently
\begin{align}
 m_\omega(s)&\le \frac1s,
 &
 -m_\omega'(s)&\ge\mathbb E[B e^{-sB}]
                  =\frac1{(1+s)^2},
 \label{van:first-child}\\
 \mathbb E\left[\left(\int_0^\infty e^{-\theta t}\xi(dt)\right)^2\right]
 &\le \frac1{\theta^2}+\frac1{2\theta},
 &&\theta>0.
 \label{van:second-moment}
\end{align}
There is positive probability of reaching zero and spending positive
time there, for every finite $\omega$.  Hence $m_\omega(1)<1$.
Since $F(0)=1-e^{-1}$, formula \eqref{van:m} also shows
$m_\omega(s)\to\infty$ as $s\downarrow0$.  Strict decrease and continuity
now give the unique root in $(0,1)$.  Algebraic rearrangement gives
\eqref{van:root}; its denominator
$\omega+(1-\omega)s=s+\omega(1-s)$ is positive on that interval.

\subsection{Branching growth and arbitrary finite seeds}

\begin{lemma}[Ordinary-family growth and a continuous-time mean bound]
\label{van:ordinary}
Put $\alpha=\alpha_\omega$ and
$\beta=-m_\omega'(\alpha)\in(0,\infty)$.
For one ordinary family started with a core of leaf count one,
\begin{equation}
 e^{-\alpha t}C^{(1)}(t)\longrightarrow
 A_1=\frac{W}{\alpha\beta}\in(0,\infty)
 \quad\text{almost surely}.
 \label{van:ordinary-limit}
\end{equation}
Here $W$ is the limiting intrinsic CMJ martingale.
For every $\sigma>\alpha$,
\begin{equation}
 \mathbb E C^{(1)}(t)
 \le \frac{e^{\sigma t}}{1-m_\omega(\sigma)}.
 \label{van:mean-bound}
\end{equation}
\end{lemma}

\begin{proof}
We verify the conditions for the CMJ law of large numbers
\cite[Section~2.2, equation~(2.13) and its almost-sure
upgrade]{IKM2024}.
Reproduction processes are independent copies by the exact
construction above.  Their mean measure is nonlattice, since the
first birth age is exponentially distributed.  The Malthusian
equation is $m_\omega(\alpha)=1$, and
\[
 0<\frac1{(1+\alpha)^2}\le\beta\le
 \int_0^\infty t e^{-\alpha t}\,dt=\alpha^{-2}<\infty.
\]
The moment bound \eqref{van:second-moment} holds for every
$\theta>0$, including a $\theta<\alpha/2$.
It supplies both the discounted $x\log x$ condition
(see \cite{Olofsson1998}) and the stronger
moment hypothesis in \cite[Remark~2.1]{IKM2024}.  The counting
characteristic $\phi(t)=\mathbf1_{\{t\ge0\}}$ is bounded and has the
required exponentially weighted regularity and summability.
For this characteristic, the weighted local suprema are bounded by
a constant times $e^{-\alpha t}\mathbf1_{\{t\ge-1\}}$, so the
characteristic condition in the almost-sure upgrade is satisfied.
Nerman's theorem gives the limit
$\beta^{-1}\int_0^\infty e^{-\alpha t}\,dt\;W
=W/(\alpha\beta)$ almost surely.

We give an explicit nondegeneracy check that also permits infinitely
many lifetime offspring. Let $\sigma_v$ be the birth time of an
individual $v$, and set
\[
 W_j=\sum_{|v|=j}e^{-\alpha\sigma_v},\qquad
 Y=\int_0^\infty e^{-\alpha t}\xi(dt).
\]
The generation martingale $(W_j)$ has mean one and the same limit $W$
as the intrinsic CMJ martingale
\cite[Section~2.2, equation~(2.12)]{IKM2024}. Independence of the child
families and monotone truncation of the discounted offspring sum give
\[
 \operatorname{Var}(W_{j+1})
 =m_\omega(2\alpha)\operatorname{Var}(W_j)+\operatorname{Var}(Y).
\]
Here $m_\omega(2\alpha)<1$ and $Y$ has finite second moment by
\eqref{van:second-moment}. Thus $W_j$ is bounded in $L^2$, so
$\mathbb E W=1$. The branching identity for $W$ implies
$q=\mathbb E[q^K]$, where $q=\mathbb P(W=0)$ and $K$ is the
possibly infinite number of children of one individual. We have
$K\ge1$ almost surely. Moreover $\mathbb P(K\ge2)\ge1/2$, since
an upward first move (probability $1/2$) leaves the core active
after its first child is born. As $q<1$, the identity forces $q=0$.
This proves $W>0$ almost surely. Local finiteness holds because the
total core-birth rate is at most the core count, giving a Yule upper
coupling.

For the mean bound, set $M(t)=\mathbb E C^{(1)}(t)$.
The ordinary branching decomposition gives
\[
 M(t)=1+\int_0^t M(t-u)\nu_\omega(du).
\]
After multiplication by $e^{-\sigma t}$ the renewal kernel has
total mass $m_\omega(\sigma)<1$, while the forcing is at most one.
Its nonnegative renewal series is bounded by
$\sum_{j\ge0}m_\omega(\sigma)^j$, which proves
\eqref{van:mean-bound}.
\end{proof}

\begin{lemma}[Positive-event inheritance of the amplitude]
\label{van:inheritance}
Let $C^{(\ell)}(t)$ be the size of the family of one initial core
whose leaf count is $\ell$.  For every integer $\ell\ge1$,
$e^{-\alpha t}C^{(\ell)}(t)$ has a finite strictly positive limit
almost surely.  The same holds for $\ell=0$ when $\omega>0$.
For $\ell=0$ and $\omega=0$ the limit is zero.
\end{lemma}

\begin{proof}
The case $\ell=1$ is Lemma~\ref{van:ordinary}.
For $\ell\ge2$, start with one ordinary core and stop at
\[
 \tau=\min\{\text{first hit of leaf count }\ell,\,
                  \text{first core birth}\}.
\]
The event $E$ that the hit occurs first has probability
\[
 \mathbb P(E)=\prod_{j=1}^{\ell-1}\frac{j}{j+1}=\frac1\ell>0.
\]
On $E$, there is still exactly one core, with leaf count $\ell$,
and no descendants.  The strong Markov future, conditionally on
the history up to $\tau$, is exactly the family started at $\ell$.
The property that $e^{-\alpha t}C(t)$ has a finite strictly positive
limit is preserved by every finite time shift.  It holds almost surely
for the ordinary family, hence also conditionally on $E$.
The Markov property therefore proves that it holds with probability
one for a family started at $\ell$.

If $\ell=0$ and $\omega>0$, an independent finite
$\operatorname{Exp}(\omega)$ activation delay is followed by the
ordinary family.  Its limiting amplitude is the ordinary amplitude
multiplied by the strictly positive discount of this delay.
For $\ell=0$ and $\omega=0$, the family consists of its one inert
ancestor forever.
\end{proof}

For a finite tree with at least three vertices, assign each initial
core its own genealogical family.  Pre-existing core edges do not enter
any rate in \eqref{van:life}, so these are finitely many independent
families with their respective initial leaf counts.  All new cores
belong to exactly one such family.  At least one initial core has
positive leaf count.  Lemma~\ref{van:inheritance} consequently gives
\begin{equation}
 e^{-\alpha t}C(t)\longrightarrow A\in(0,\infty)
 \quad\text{almost surely}.
 \label{van:seed-amplitude}
\end{equation}

We will also need a mean bound for this seed.  The reproduction
intensity of each exceptional initial core is at most one.
Every child begins an ordinary family, so \eqref{van:mean-bound}
and a convolution with this initial reproduction measure show that,
for every $\sigma>\alpha$,
\begin{equation}
 \mathbb E C(t)\le K_\sigma e^{\sigma t}
 \label{van:seed-mean}
\end{equation}
for some finite constant depending on the seed.  For example, if
there are $k$ initial cores, one may take
$K_\sigma=k\bigl(1+[\sigma(1-m_\omega(\sigma))]^{-1}\bigr)$.
This is an auxiliary continuous-time bound; no assertion about
discrete-time mean asymptotics is being used.

\subsection{The exact population clock}

A core birth removes one leaf from its parent and supplies one leaf
to the newborn core.  It leaves the \emph{total} leaf count unchanged.
All other vertex births increase that count by one.  Hence
\begin{equation}
 L(t)\ \text{jumps by }+1\text{ at rate }\ L(t)+\omega C_0(t).
 \label{van:leaf-rate}
\end{equation}
Choose $\alpha<\sigma<1$ in \eqref{van:seed-mean}.
The compensated counting process yields
\begin{equation}
 e^{-t}L(t)=L(0)+\omega\int_0^t e^{-u}C_0(u)\,du+\mathcal M_t,
 \label{van:leaf-martingale}
\end{equation}
where $\mathcal M$ is a square-integrable martingale on every finite
interval.  The drift integral has finite expectation as $t\to\infty$.
Moreover the mean version of \eqref{van:leaf-martingale} gives
\[
 \mathbb E L(t)
 =e^t\left(L(0)+\omega\int_0^t e^{-u}\mathbb E C_0(u)\,du\right)
 \le K e^t
\]
for a finite $K$.  It follows that
\[
 \mathbb E\langle\mathcal M\rangle_\infty
 =\int_0^\infty e^{-2u}
       \bigl(\mathbb E L(u)+\omega\mathbb E C_0(u)\bigr)\,du<\infty.
\]
Thus both terms on the right of \eqref{van:leaf-martingale} converge
almost surely and $e^{-t}L(t)\to V<\infty$.
From \eqref{van:leaf-rate}, $L$ admits a pathwise lower coupling
with a rate-one Yule process started with $L(0)$ particles.
Its $e^{-t}$-normalized limit is strictly positive almost surely.
Therefore $V>0$.

Since $\alpha<1$, \eqref{van:seed-amplitude} implies
$e^{-t}C(t)\to0$, whence $e^{-t}N(t)\to V$.
For the hitting time $\tau_n$ of population $n$, nonexplosion and
the Yule lower bound imply $\tau_n\to\infty$.
Consequently
\[
 \frac{C_n}{n^\alpha}
 =\frac{e^{-\alpha\tau_n}C(\tau_n)}
        {(e^{-\tau_n}N(\tau_n))^\alpha}
 \longrightarrow A V^{-\alpha}\in(0,\infty)
 \quad\text{almost surely}.
\]
A two-vertex seed is covered by applying the argument after its
first event and absorbing the finite time shift into $A$ and $V$.

\subsection{Effective computation of the root}

The transform admits the explicit convergent series
\begin{equation}
 F(s)=e^{-1}\sum_{j=0}^\infty\frac1{j!(s+j+1)},
 \qquad s>-1.
 \label{van:series}
\end{equation}
For $0\le s\le1$, truncation after $j=J$ has error at most
\[
 \frac{2}{(J+2)(J+1)!}.
\]
This follows by bounding $s+j+1\ge J+2$ in the tail and using
$\sum_{j\ge J+1}1/j!\le2/(J+1)!$.
Rational enclosures of $e^{-1}$ and a computable $\omega$ therefore
give certified enclosures of $m_\omega(s)$ at positive rational $s$.

The endpoint limits proved above allow effective discovery of rational
$a,b$ with $0<a<\alpha_\omega<b<1$ and
$m_\omega(a)>1>m_\omega(b)$.
Strict monotonicity gives certified root isolation.  To avoid relying
on deciding equality at a trial point, use the slope bound
\eqref{van:slope}: if a computed enclosure proves
$|m_\omega(c)-1|\le\delta$ for $c\in(0,1)$, then
$|\alpha_\omega-c|\le4\delta$.
Interval evaluation and this estimate supply terminating rational
enclosures of any requested width.

\begin{proof}[Completion of the proof of Theorem~\ref{van:main}]
The excursion calculation and renewal identity establish
\eqref{van:m} and \eqref{van:root}.  The first-child and endpoint
estimates prove uniqueness, the location $(0,1)$ and
\eqref{van:slope}.  The CMJ theorem, positive-event inheritance,
finite-family decomposition and leaf-clock argument give all
almost-sure limits in \eqref{van:limits}.  The last subsection
proves effectivity.
\end{proof}

\begin{remark}[Model distinction]
\label{van:scope}
For the original isotropic-redirection friend tree, a core may
redirect to another core.  Neighbouring environments then influence
its future leaf count, and the independent reproduction construction
\eqref{van:life} no longer holds.  The theorem supplies an exact
characterization for the modified VAN family only.  The original friend-tree exponent is not identified by this
characteristic equation. We use no discrete-time $L^1$ or
expected-amplitude conclusion.
\end{remark}


\section{Irrationality by arithmetic specialization}
\label{vana:section}

The modified VAN models admit an arithmetic specialization argument
that is not presently available for the original isotropic-redirection
model.  The key distinction is that their exact characteristic equation
uses a single explicitly controlled transform,
\begin{equation}\label{vana:F}
 F(s)=\int_0^1 x^s e^{x-1}\,dx,\qquad s>-1.
\end{equation}
We first prove the required specialization theorem directly from the
classical Siegel--Shidlovskii theorem.

\begin{theorem}[Transcendence of rational specializations]
\label{vana:rational-specialization}
For every rational number $r>-1$, the number $F(r)$ is transcendental.
\end{theorem}

\begin{proof}
Put $\lambda=r+1>0$ and define
\begin{equation}\label{vana:f}
 f_\lambda(z)=\sum_{n=0}^\infty\frac{z^n}{(\lambda+1)_n}
             =\sum_{n=0}^\infty c_n\frac{z^n}{n!},
 \qquad c_n=\frac{n!}{(\lambda+1)_n},
\end{equation}
where $(u)_0=1$ and $(u)_n=u(u+1)\cdots(u+n-1)$.
The ratio of successive terms proves that $f_\lambda$ is entire.
Substituting $u=1-x$ in~\eqref{vana:F} and integrating the exponential
series gives
\begin{align}
 F(r)
 &=\sum_{n=0}^\infty\frac{(-1)^n}{n!}
      \int_0^1u^n(1-u)^{\lambda-1}\,du\notag\\
 &=\frac1\lambda
      \sum_{n=0}^\infty\frac{(-1)^n}{(\lambda+1)_n}
  =\frac{f_\lambda(-1)}{\lambda}.
 \label{vana:beta-identity}
\end{align}
The interchange is justified by the uniformly absolutely convergent
exponential series and the integrable weight
$(1-u)^{\lambda-1}$.

We verify the arithmetic conditions for $f_\lambda$ to be an
$E$-function.  Write $\lambda=a/b$ with relatively prime positive
integers $a,b$.  Then
\begin{equation}\label{vana:coefficients}
 c_k=\frac{b^k k!}{\prod_{j=1}^k(a+bj)}.
\end{equation}
We claim that the reduced denominators of all $c_0,\ldots,c_n$ divide
the single integer
\begin{equation}\label{vana:denominator}
 D_n=\operatorname{lcm}(1,2,\ldots,a+bn).
\end{equation}
Indeed, if a prime $p$ divides $b$, it divides none of the factors
$a+bj$, so it contributes no denominator.  If $p\nmid b$, the
congruence $a+bj\equiv0\pmod{p^h}$ specifies one residue class
modulo $p^h$.  Hence, for every $k\le n$,
\[
 \#\{1\le j\le k:p^h\mid a+bj\}
 \le \left\lfloor\frac{k}{p^h}\right\rfloor+1.
\]
There are no such factors when $p^h>a+bk$.  Comparing the valuation
of the denominator in~\eqref{vana:coefficients} with Legendre's
formula for $v_p(k!)$ gives
\[
 \max\{0,-v_p(c_k)\}
 \le \left\lfloor\log_p(a+bk)\right\rfloor
 \le \left\lfloor\log_p(a+bn)\right\rfloor
 =v_p(D_n).
\]
This proves the common-denominator assertion.
The elementary bound
$\log\operatorname{lcm}(1,\ldots,M)=O(M)$ therefore yields
$\log D_n=O_\lambda(n+1)$.  Also
\[
 0<c_k=\prod_{j=1}^k\frac{j}{\lambda+j}\le1.
\]
The coefficients are rational, so the same bound holds for all their
algebraic conjugates.  Thus the logarithmic height of
$(c_0,\ldots,c_n)$ is $O_\lambda(n+1)$, as required in the
$E$-function definition.

The coefficient recurrence yields
\begin{equation}\label{vana:first-order}
 z f_\lambda'(z)=(z-\lambda)f_\lambda(z)+\lambda.
\end{equation}
In particular, differentiation gives the homogeneous equation
\[
 z f_\lambda''(z)+(\lambda+1-z)f_\lambda'(z)-f_\lambda(z)=0,
\]
with coefficients in $\mathbb Q[z]$.  This completes the verification
that $f_\lambda$ is an $E$-function.

Every coefficient in~\eqref{vana:f} is nonzero, so $f_\lambda$ is
not a polynomial.  An entire function algebraic over $\mathbb C(z)$
must be a polynomial: a polynomial relation and the usual bound for
roots of a polynomial give polynomial growth outside a disk, after
which Cauchy's estimates force all sufficiently high Taylor
coefficients to vanish.  Consequently $f_\lambda$ is transcendental
over $\mathbb C(z)$.

The two $E$-functions $1,f_\lambda$ satisfy the homogeneous system
\begin{equation}\label{vana:system}
 \frac{d}{dz}
 \begin{pmatrix}1\\f_\lambda(z)\end{pmatrix}
 =
 \begin{pmatrix}
  0&0\\
  \lambda/z&1-\lambda/z
 \end{pmatrix}
 \begin{pmatrix}1\\f_\lambda(z)\end{pmatrix}.
\end{equation}
Its coefficient matrix has common denominator $T(z)=z$.
The Siegel--Shidlovskii theorem in the form of
\cite[Theorem~1.1]{Beukers2006} applies at the nonzero algebraic
point $\xi=-1$, since $\xi T(\xi)=1$.  It gives
\[
 \operatorname{trdeg}_{\overline{\mathbb Q}}
    \overline{\mathbb Q}\bigl(1,f_\lambda(-1)\bigr)
 =
 \operatorname{trdeg}_{\overline{\mathbb Q}(z)}
    \overline{\mathbb Q}(z)\bigl(1,f_\lambda(z)\bigr)
 =1.
\]
Thus $f_\lambda(-1)$ is transcendental.  Equation
\eqref{vana:beta-identity} proves the assertion because
$\lambda$ is a nonzero rational number.
\end{proof}

\begin{theorem}[Irrationality of finite algebraic-parameter VAN roots]
\label{vana:irrationality}
Let $\omega\in\overline{\mathbb Q}\cap[0,\infty)$ be finite.
The unique root $\alpha_\omega\in(0,1)$ of
\begin{equation}\label{vana:root}
 F(\alpha_\omega)
 =\frac{(1-\alpha_\omega)(\alpha_\omega+\omega)}
        {\omega+(1-\omega)\alpha_\omega}
\end{equation}
is irrational.  By Theorem~\ref{van:main}, this is the almost-sure core-growth
exponent of $\operatorname{VAN}(\omega)$.  In
particular, the PAN and DAN exponents are irrational.
\end{theorem}

\begin{proof}
Theorem~\ref{van:main} supplies the existence and uniqueness of
this root in $(0,1)$.  The arithmetic assertion needs
only the root equation.  If $\alpha_\omega$ were rational, its
right-hand side would be algebraic: its denominator satisfies
\[
 \omega+(1-\omega)\alpha_\omega
 =\alpha_\omega+\omega(1-\alpha_\omega)>0.
\]
But Theorem~\ref{vana:rational-specialization} makes its left-hand
side transcendental.  This is a contradiction.  The cases
$\omega=0$ and $\omega=1$ give PAN and DAN, respectively.
\end{proof}

\paragraph{Connection with polynomial exponential integrals.}
For a rational candidate $r=p/q\in(0,1)$ in lowest terms, integration
by parts followed by $x=t^q$ gives the useful normal form
\begin{equation}\label{vana:EIT-form}
 F(p/q)
 =1-\frac pq e^{-1}\int_0^1 x^{p/q-1}e^x\,dx
 =1-p e^{-1}\int_0^1 t^{p-1}e^{t^q}\,dt.
\end{equation}
Thus rational specialization places the problem in the polynomial
exponential-integral setting.  The polynomial amplitude
$p t^{p-1}$ already has degree at most $q-2$.  It is a nonzero
reduced amplitude for the twisted derivative
\[
 \mathcal D_q R=R'+q t^{q-1}R,
 \qquad
 \frac{d}{dt}\bigl(R(t)e^{t^q}\bigr)
       =(\mathcal D_qR)(t)e^{t^q}.
\]
Indeed, for a nonzero polynomial $R$, the highest-degree term of
$\mathcal D_qR$ has degree $\deg R+q-1$, so no nonzero polynomial
of degree at most $q-2$ can be of this form.  Allowing rational $R$
does not create an exception: a finite pole of $R$ gives an
uncancelled pole of one higher order in $R'$, whereas
$q t^{q-1}R$ has no pole of that higher order.  A rational
$R$ with polynomial $\mathcal D_qR$ is therefore polynomial.
This supplies a concrete reduced exponential-integral class to
which an arithmetic specialization theorem can be applied.
For this family, Theorem~\ref{vana:rational-specialization} supplies
that theorem through classical $E$-function theory.

\paragraph{Limits of the conclusion.}
The argument uses rationality of the candidate exponent to obtain an
$E$-function with controlled arithmetic coefficients.  It does not
exclude algebraic irrational roots and therefore does not prove that
$\alpha_\omega$ is transcendental.  The singular model
$\omega=\infty$, whose limiting exponent is $1$, is outside the
theorem.  These are modified local redirection models: no equation
\eqref{vana:root}, and no irrationality or transcendence conclusion,
is asserted for the original isotropic-redirection exponent $\mu$.

\section{A transfer criterion and its scope}
\label{vana:transfer-section}

\begin{proposition}[Transfer with specialization nonvanishing]
\label{vana:transfer}
Let $I\subset(-1,\infty)$ be a real interval.  Suppose an exact
characterization of a target number $\alpha\in I$ gives
\begin{equation}\label{vana:determinant-form}
 D(\alpha)=0,\qquad
 D(s)=\frac{P(s,F(s))}{Q(s,F(s))},
 \qquad P,Q\in\overline{\mathbb Q}[s,X],
\end{equation}
where the denominator is nonzero on $I$.  If
\begin{equation}\label{vana:specialization-nonzero}
 P(r,X)\not\equiv0
 \quad\text{for every }r\in I\cap\mathbb Q,
\end{equation}
then $\alpha$ is irrational.
\end{proposition}

\begin{proof}
If $\alpha=r\in\mathbb Q$, then
\eqref{vana:determinant-form} gives $P(r,F(r))=0$.
Condition~\eqref{vana:specialization-nonzero} makes this a nonzero
polynomial relation over $\overline{\mathbb Q}$ for $F(r)$,
contradicting Theorem~\ref{vana:rational-specialization}.
\end{proof}

The proposition applies, for example, when an exact finite system
has a characteristic determinant of the displayed form.  It is a
criterion to verify, not a claim that the original friend-tree
kernel has such a reduction.  The specialization condition is
essential.  A nonzero bivariate polynomial can vanish identically
as a polynomial in $X$ at a rational parameter: for example,
$P(s,X)=(s-r_0)(X+1)$ vanishes at $s=r_0$ for every value of
$F(r_0)$.  If $P(s,X)=\sum_jP_j(s)X^j$, the needed condition is
precisely that the coefficient polynomials $P_j$ have no common
rational zero in $I$.  This is the nonvanishing step that an exact
elimination must preserve.


\section{Certified numerical examples}
\label{comp:section}

The factorial-tail bound in~\eqref{van:series} and rational alternating
bounds for $e^{-1}$ permit root isolation entirely in rational
arithmetic. For fixed $s>0$ and $\omega\ge0$, the expression
\[
 \frac{(s+\omega)(1-f)}{s(s+\omega-\omega f)}
\]
is strictly decreasing in $f$ wherever its denominator is positive.
Thus an interval enclosure of $F(s)$ gives an enclosure of
$m_\omega(s)$ with the endpoint order reversed. The two strict signs
$m_\omega(a)>1>m_\omega(b)$ certify $a<\alpha_\omega<b$.

The script \texttt{scripts/check\_van\_root.py} implements this
calculation using Python's exact rational arithmetic. Its recorded
intervals have width $2^{-90}$. Table~\ref{comp:table} displays coarser
decimal intervals of width $10^{-24}$ containing those exact intervals.
The full rational endpoints, series truncation order, and positive
rational sign margins are recorded in
\texttt{verification/van\_root\_checks.json}.

\begin{table}[htbp]
\centering
\small
\begin{tabular}{@{}ccl@{}}
\toprule
Model & $\omega$ & Certified enclosure of $\alpha_\omega$ \\
\midrule
PAN & $0$ & $(0.549734972504018815207626,\;0.549734972504018815207627)$ \\
DAN & $1$ & $(0.771192302407000256316924,\;0.771192302407000256316925)$ \\
VAN$(2)$ & $2$ & $(0.839188216788971794281543,\;0.839188216788971794281544)$ \\
\bottomrule
\end{tabular}
\caption{Rationally certified examples for the local models.}
\label{comp:table}
\end{table}

The script also checks the common-denominator divisibility in
\eqref{vana:denominator} on a finite collection of rational parameters.
Those checks support reproducibility of the algebra; the
all-parameter denominator proof, the branching identification, and
the irrationality theorem are established analytically in the text.

\section*{Acknowledgments and preparation}
The research questions arose from the study of the original random
friend-tree exponent. AI tools assisted with literature searches,
proof development, exact-arithmetic code, and manuscript preparation.
The mathematical arguments and their stated scope are presented here
for independent review.

\begin{thebibliography}{9}

\bibitem{Beukers2006}
F.~Beukers,
\emph{A refined version of the Siegel--Shidlovskii theorem},
Annals of Mathematics \textbf{163} (2006), 369--379.
\href{https://annals.math.princeton.edu/wp-content/uploads/annals-v163-n1-p08.pdf}{Publisher PDF}.

\bibitem{HRKZ2026}
H.~Hartle, P.~L.~Krapivsky, S.~Redner, and Y.~Zhang,
\emph{Anomalous scaling in redirection networks},
Physical Review E \textbf{114} (2026), 034314.
\href{https://doi.org/10.1103/9lmp-pydp}{doi:10.1103/9lmp-pydp};
\href{https://arxiv.org/abs/2604.01540}{arXiv:2604.01540}.

\bibitem{IKM2024}
A.~Iksanov, K.~Kolesko, and M.~Meiners,
\emph{Asymptotic fluctuations in supercritical Crump--Mode--Jagers processes},
\href{https://arxiv.org/abs/2109.00867v3}{arXiv:2109.00867v3} (2024).
Section~2.2, especially Remark~2.1, equation~(2.13), and its
almost-sure upgrade.

\bibitem{KR2017}
P.~L.~Krapivsky and S.~Redner,
\emph{Emergent Network Modularity},
Journal of Statistical Mechanics: Theory and Experiment (2017), 073405.
\href{https://doi.org/10.1088/1742-5468/aa7a3f}{doi:10.1088/1742-5468/aa7a3f};
\href{https://arxiv.org/abs/1706.01514}{arXiv:1706.01514}.

\bibitem{Olofsson1998}
P.~Olofsson,
\emph{The $x\log x$ condition for general branching processes},
Journal of Applied Probability \textbf{35} (1998), 537--544.
\href{https://www.stat.rice.edu/~olofsson/xlogx.pdf}{Author PDF}.

\end{thebibliography}
\end{document}
